% Conclusion
\section{Conclusion}
\label{sec:conclusion}

We introduced SPOQ (Specialist Orchestrated Queuing), a methodology for
multi-agent software engineering that addresses coordination, quality,
and human oversight challenges in prior approaches.

\subsection{Summary of Contributions}

\paragraph{Wave-Based Topological Dispatch.}
SPOQ computes parallel execution waves from task dependency graphs.
Experiment~1 demonstrates the scheduler's behavior under two regimes:
on unbounded synthetic DAGs, wave dispatch reaches the critical-path
lower bound (ratio 1.03--1.11) with speedups up to 14.3$\times$,
showing that the algorithm captures essentially all available
parallelism; on a 2-slot local backend running real LLM calls, it
delivers a stable 1.4$\times$ speedup that matches the hardware
concurrency ceiling, showing that the result survives realistic
backend constraints. Field deployments report 1.3--5.3$\times$
wall-clock gains across diverse production epics.

\paragraph{Dual Validation Gates.}
By validating both planning quality (10 metrics, 95/90 threshold) and
code quality (10 metrics, 95/80 threshold), SPOQ catches issues at the
stages where they are cheapest to fix. Experiment~3 shows that dual
validation reduces defects per task from 0.34 to 0.20, increases test
pass rate from 91.25\% to 99.75\%, and eliminates static warnings
across the benchmark suite.

\paragraph{Human-as-an-Agent Integration.}
The HaaA paradigm positions the human specialist as a high-value agent
within the orchestration loop, enabling bidirectional collaboration that
amplifies both human judgment and agent productivity. Experiment~4
demonstrates that human-assisted planning further reduces residual
defects from 0.47 to 0.03 per task and lifts test pass rate from 96.5\%
to 99.75\%, with a parallel improvement in plan quality observed under a
local open-weights provider.

\paragraph{Three-Tier Agent Hierarchy.}
SPOQ's Opus/Sonnet/Haiku hierarchy optimizes cost-quality tradeoffs by
matching agent capabilities to role requirements: high capability for
execution, balanced for review, fast-cheap for triage.

\paragraph{Cross-Provider Replication.}
All four experiments were replicated against a locally hosted
Qwen3.6-35B-A3B model served by \texttt{llama.cpp}. The direction and
significance of the SPOQ gains are preserved across provider families:
structured planning recovers 35 points of coverage and 52.5 points of
parallelism that the unaided Qwen baseline loses, eliminates the
cyclic-plan failure mode (2/4 $\to$ 0/4), and the validation ablation
maintains the same monotonic Full-SPOQ $>$ Code-Validation $>$
No-Validation ordering observed under Claude. This supports the claim
that the improvements stem from orchestration rather than from any
specific model's capabilities.

\paragraph{Practical Lessons.}
Through controlled benchmarks and a 17-repository deployment study we
identified and addressed failure modes including runaway retry loops,
lock file contention, and context window exhaustion.

\subsection{Vision for AI-Native Engineering}

SPOQ represents early steps toward a future where engineers become
architects and validators rather than line-by-line implementers, where
quality is built into process through structured gates, and where human
expertise scales through orchestration. Much work remains to realize this
vision fully. We hope SPOQ provides a useful framework for researchers
and practitioners exploring the frontier of AI-assisted software
development.

\subsection*{Acknowledgments}

We thank the Claude Code development team for the tooling infrastructure
that enabled this research. We also thank Ross Sylvester (CEO, Adrata) and
John Armbruster (Founding Engineer, Notary Everyday), whose adoption of
SPOQ and candid feedback helped refine the methodology.
